\section{Related Work}
\label{sec:related}

\para{Order sensitivity in in-context learning.}
In-context learning with demonstrations was popularized by GPT-3~\cite{brown2020language}.
\citet{lu2022fantastically} showed that permuting the same demonstrations can move GPT-2 and GPT-3 classification accuracy from near chance to near state of the art, that this sensitivity did not shrink with model size, and that good orders do not transfer across models; they proposed selecting orders with entropy statistics on a synthetic probing set.
\citet{zhao2021calibrate} traced part of this instability to majority-label, recency and common-token biases, and corrected them with contextual calibration.
Unlike these studies, which used base models of the GPT-2/3 era and mostly reported ranges over sampled permutations, we study a modern instruction-tuned family, use a sensitivity metric that is unbiased in the number of sampled orders, and test whether order effects exceed item-level and numerical noise.

\para{Position effects in long contexts.}
\citet{liu2024lost} found that LLMs use information at the start and end of long contexts better than information in the middle.
That work studies retrieval of one relevant passage among many distractors.
We instead study short prompts in which every demonstration is relevant, and ask how much each position steers the output.
We measure this with direct interventions on one position at a time rather than by moving a single relevant item.

\para{What demonstrations contribute.}
\citet{min2022rethinking} found that the format and label space of the demonstrations matter more than whether their labels are correct.
Our label-flip probe builds on this observation: it asks whether flipping a single label, at a given position, moves the model toward that label.
We find that it often does, but that for mid-size models on NLI a flipped final label moves predictions \emph{away} from it.
Work on demonstration \emph{selection}~\cite{liu2022makes,rubin2022learning} chooses which examples to show; we hold the selected set fixed and vary only the order, so our results describe the variance that remains after selection.

\para{Evaluation variance.}
Prompt-format choices can change LLM accuracy by large margins~\cite{sclar2024quantifying}, which suggests that many reported gains are within the noise of prompt design.
We contribute two checks for the order dimension of this problem: a split-half reliability test that asks whether a good order stays good on held-out test items, and an identical-prompt rerun that measures the non-determinism of batched bf16 inference in vLLM~\cite{kwon2023efficient}.
